The number $1066$ immediately reminds us of the battle of Hastings and the victory of William the Conqueror, Duke of Normandy;
$1776$ of the independence of the United States; $1789$ of the French
Revolution. When we read $3.14159$ we know that it is an approximate value of $\pi$
(written, as in English, with a decimal point, where French uses a comma).
On the other hand, when we read the digit $2$, we think of nothing in particular and we are almost certain
to wonder
``two what?''.
All this is to say that when we see a number in a
program it is not always easy to understand immediately what it is about. To make it clear,
we can, in \python, name the number with a constant declaration placed at the
beginning of the program:

{
\parindent \parindentprog
\medskip

\vspace{\neglskip}\mbox{\blc {\idt frenchRevolution\ =\ 1789\ \ }{\cmt \#\ Date\ of\ the\ revolution}}\par

\medskip }

\noindent and then use the name {\idt frenchRevolution} instead of the number
{\small\rm 1789} in the rest of the program.


\begin{exemple}[Maltese cross]{\rm
In the program below, which draws a Maltese cross, we
use a constant {\idt length} to define the length of an
arm (that is to say the distance from its base to the centre of the cross) and
another one, {\idt width}, to define the width of an arm
(that is to say the distance between the ends of its base):\hfill\\[22mm]
\includePostScript{CroixDeMalte_1}{0.95}{0.95}%
\vspace*{-22mm}%
\input python/MalteseCross_1.tex

To choose the best possible shape for
this cross, there is only one place in the program where these constants
need to be changed. For example, to obtain the smallest cross shown in
figure \ref{fig:crmadiff} below, it is enough to take:\par
{
\parindent \parindentprog
\medskip

\vspace{\neglskip}\mbox{\blc {\idt width\ =\ 10\ \ }{\cmt \#\ Width\ of\ an\ arm}}\par
\vspace{\neglskip}\mbox{\blc {\idt length\ =\ 15\ \ }{\cmt \#\ Length\ of\ an\ arm}}\par

\medskip }
\begin{figure}[htbp]
 \includevspacePostScript[2.5cm]{CroixDeMalteDiff}{0.95}{0.95}
 \caption{Maltese crosses $10\times 15$, $15\times 20$, $20\times 25$, $25\times 30$,
 $10\times 35$, $35\times 35$}
 \label{fig:crmadiff}
\end{figure}
\noindent $\Box$ }\end{exemple}

Note that the numbers {\idt 20} and {\idt 30} (or {\idt 10} and {\idt 15})
are {\em integer} numbers. They are used as the lengths of the sides of the
squares of the grid in the commands {\idt gsX} and {\idt gsY} which, like the command
{\idt gs}, also accept numbers with a decimal part, written with a point, such as {\idt 20.5},
which are called {\em floating-point} numbers. So we could just as well have written
`{\idt width\ =\ 20.5}'. Integers (such as {\idt 20}) and floating-point numbers
(such as {\idt 20.5} or {\idt 20.0}) are two different kinds of numbers for \python,
but it mixes them without difficulty in computations. We will see
why they are different in chapter 19 on rational expressions.

We indicate that we use one or sev\-er\-al {\em named constants}
by {\em constant declarations}:\index{Constant declaration}

\begin{center}
\begin{tabular}{@{}|p{12.5cm}|@{}} \hline
{\em Constant declaration:}\\ \hline \hline
\begin{minipage}[t]{12.5cm}
{\medskip
\mbox{\blc\ \ \ \ \idt {\small\rm identifier1}\ $=$\ {\small\rm number1}}\par
\mbox{\blc\ \ \ \ \idt {\small\rm identifier2}\ $=$\ {\small\rm number2}}\par
\mbox{\blc\ \ \ \ \small\rm  \ldots}\par
\medskip}
\end{minipage}\\ \hline
\end{tabular}
\end{center}

\noindent Each constant declaration is written on one line, at the beginning of the line
(without indentation), like `$x = 3$' in mathematics. It is customary to put a space
\hspace{0.3em}\blanc\hspace{0.3em} before and after the $=$ sign, but it is not compulsory.
As for procedure names, the identifier must follow the rules for
identifiers of \python. The custom is to start it with a lower-case
letter. Upper-case and lower-case letters are distinguished: if we declare
`{\idt Width\ =\ 20}' and then use {\idt width}, \python\ answers
`{\tt NameError: name 'width' is not defined. Did you mean: 'Width'?}'.
After these declarations, we can use the constant names
`{\small\rm identifier1}' instead of `{\small\rm number1}' and `{\small\rm identifier2}'
instead of `{\small\rm number2}', etc.

To tell the truth, \python\ has no real constants: nothing prevents us from giving, further on
in the program, another value to the name {\idt width} (we will see in chapter 22 that this is
how {\em variables} are used). We speak of a constant because the program
gives a value to this name only once, at the beginning, and never changes it afterwards.

In a \python\ program, a declaration defines constants (with the $=$ sign) or
procedures (with the word {\mcl def}). We will see later that it can also define functions and
variables:\index{Declaration}

\begin{center}
\begin{tabular}{@{}|p{12.5cm}|@{}} \hline
{\em Declaration:}\\ \hline \hline
\begin{minipage}[t]{12.5cm}
{\medskip
\mbox{\blc\ \ \ \ \small\rm  Constant declaration}\par
\mbox{\blc\ \ \tiny\rm  or}\par
\mbox{\blc\ \ \ \ \small\rm  Procedure declaration}\par
\medskip}
\end{minipage}\\ \hline
\end{tabular}
\end{center}

Constant or procedure declarations are placed, in any
order, before the instructions of the program:\index{Program}

\begin{center}
\begin{tabular}{@{}|p{12.5cm}|@{}} \hline
{\em \python\ program:}\\ \hline \hline
\begin{minipage}[t]{12.5cm}
{\medskip
\mbox{\blc\ \ \ {\mcl from}\hspace{0.3em}\blanc\hspace{0.3em}{\idt LibRobot}\hspace{0.3em}\blanc\hspace{0.3em}{\mcl import}\ {\idt *}}\par
\mbox{\blc\ \ \ \cmt \#\ \small\rm Comment explaining what the program does}\par
\mbox{\blc\ \ \ \small\rm  \ldots\ Declarations of the program \ldots}\par
\mbox{\blc\ \ \ \small\rm  \ldots\ Instructions of the program\footnote{\ Using the declared constants
and procedures.} (one per line or separated by `$;$') \ldots}\par
\medskip}
\end{minipage}\\ \hline
\end{tabular}
\end{center}

The main rule to follow is that every identifier must be declared before it is
used (otherwise, when the \robot\ reaches the instruction that uses it, \python\ stops with
a {\tt NameError} error).

\needspace{6\baselineskip}\noindent \rule{137mm}{0.1mm}\par\nopagebreak %---------------------------------------------
\begin{exercice}{Write \python\ programs with named constants
to draw:

\tiret the singular propeller of figure \ref{fig:hesing} below;

\tiret the duckling of figure \ref{fig:canar} below.

\noindent $\Box$}\end{exercice}
\begin{figure}[htbp]
 \includevspacePostScript[4.6cm]{HeliceSinguliere}{0.95}{0.95}
 \caption{Singular propeller}
 \label{fig:hesing}
\end{figure}
\begin{figure}[htbp]
 \includevspacePostScript[4.2cm]{Canardeau}{0.95}{0.95}
 \caption{Duckling hatching from its egg}
 \label{fig:canar}
\end{figure}
\begin{corrige}{
We give one program; the other one is in the folder {\idt book/\allowbreak 11\_constants}.
The horizontal side of the squares of the grid of the $i$-th
leaf of a blade is {\idt xbase * i}. It is an integer that can be given
directly to the command {\idt gsX}:
\input python/SingularPropeller.tex
\noindent $\Box$}\end{corrige}
\nopagebreak

%%%%%%%%%%%%%%%%%%%%%%%%%%%%%%%%%%%%%%%%%%%%%%%%%%%%%%%%%%%%%%%%%%%%%%%%%%%
\noindent \framebox[\textwidth]{
\begin{minipage}{\myparboxwidth}
\setlength{\parindent}{2em}
\small \sf
\vspace{\abovedisplayskip}
\section*{Addenda}
Computer integers are an approximation of the \href{https://en.wikipedia.org/wiki/Integer}{integers} of mathematics (but not too large ones). Floating-point numbers (more precisely \href{https://en.wikipedia.org/wiki/Floating-point_arithmetic}{floating-point} numbers) are a subset of the \href{https://en.wikipedia.org/wiki/Rational_number}{rational numbers} (and therefore of the \href{https://en.wikipedia.org/wiki/Real_number}{real numbers}) written in \href{https://en.wikipedia.org/wiki/Floating-point_arithmetic}{floating-point} notation.

\vspace{\abovedisplayskip}
\end{minipage}
}

%%%%%%%%%%%%%%%%%%%%%%%%%%%%%%%%%%%%%%%%%%%%%%%%%%%%%%%%%%%%%%%%%%%%%%%%%%%%
%%%%%%%%%%%%%%%%%%%%%%%%%%%%%%%%%%%%%%%%%%%%%%%%%%%%%%%%%%%%%%%%%%%%%%%%%%%
\noindent \framebox[\textwidth]{
\begin{minipage}{\myparboxwidth}
\setlength{\parindent}{2em}
\small \sf
\vspace{\abovedisplayskip}
A {\em constant declaration} {\idt name\ =\ number} gives a name
to a particular number. Afterwards the name, which is more meaningful, can be
used instead of the number. The immediate advantage is that to change the number it is
enough to do so in the constant declaration, whereas otherwise one has to look for all its
occurrences in the program.

This can be compared with the notion of variable in mathematics, that is to say with the use
of letters $x$, $y$, $z$ \ldots\ to denote arbitrary numbers. The difference is that in
\python\ the name of a constant denotes a particular value whereas in mathematics it
is an arbitrary value.

Historically, the first mathematical problems dealt with by the Babylonians (1800 years BC)
included an algorithmic solution given for one example and separate numerical tables
 to solve other similar problems \cite{ritter1}.
In the Middle Ages people used rhetorical language such as `a number added to its square
makes 12' for `$x + x^{2} = 12$'. Chuquet\footnote{\ \href{https://en.wikipedia.org/wiki/Nicolas_Chuquet}{Nicolas Chuquet}, French
mathematician, c.~1445--1500.}, author of the oldest treatise on algebra written by a
Frenchman (1484), used $1^{1}$ to denote the unknown, $1^{2}$ to denote its
square and the symbols $\bar{p}$ and $\bar{m}$ for addition and subtraction.
He could write `$12\;\bar{p}\;3^{1}$ egaulx a $4^{2}$' (Middle French for ``equal to'') for `$12 + 3x = 4x^{2}$'.
Viète\footnote{\ \href{https://en.wikipedia.org/wiki/François_Viète}{François Viète}, French mathematician, 1530--1603.} took an
important step towards symbolism in algebra by introducing the use of letters to denote
unknown quantities. Around 1590 he wrote:
`$1\;$ QC $\;-\;15\;$QQ$\;+\;85\;$C$\;+\;225\;$Q$\;+\;274\;$N equal $120$', that is, if
C means cube, hence $x^{3}$, Q means square, hence $x^{2}$, and N means number, hence $x$, the equation
$x^{5} - 15x^{4} + 85x^{3} + 225x^{2} + 274x = 120$ (\cite{benoit1}). It took four French
mathematicians,
Desargues\footnote{\ \href{https://en.wikipedia.org/wiki/Girard_Desargues}{Girard Desargues}, French mathematician, 1591--1661.},
Descartes\footnote{\ \href{https://en.wikipedia.org/wiki/René_Descartes}{René Descartes}, French mathematician, 1596--1650.},
Fermat\footnote{\ \href{https://en.wikipedia.org/wiki/Pierre_de_Fermat}{Pierre de Fermat}, French mathematician, 1601--1665.} and
Pascal\footnote{\ \href{https://en.wikipedia.org/wiki/Blaise_Pascal}{Blaise Pascal}, French mathematician, 1623--1662.} (the latter having been chosen
by Wirth to name his programming language \pascal), to enter the modern
era of mathematics.

\vspace{\abovedisplayskip}
\end{minipage}
}
%%%%%%%%%%%%%%%%%%%%%%%%%%%%%%%%%%%%%%%%%%%%%%%%%%%%%%%%%%%%%%%%%%%%%%%%%%%%
